\subsection{Náročnost provozu}
\label{subsec:criteria-operational}

Náročnost provozu určuje, jaké zdroje jsou potřeba k dlouhodobému provozování modelu v produkčním prostředí systému Kelvin.
Zahrnuje finanční náklady, nároky na infrastrukturu i lidskou práci spojenou s údržbou.
Struktura nákladů se zásadně liší podle zvoleného způsobu nasazení.

U cloudového API jsou náklady průběžné a přímo úměrné objemu zpracovaných dat.
Jak bylo popsáno v \secref[sekci]{subsec:llm-tokens-parameters-context}, cenové modely poskytovatelů jsou založeny na počtu zpracovaných tokenů.

Pro odhad typické spotřeby tokenů byla analyzována sada studentských odevzdání z reálného testovacího provozu systému Kelvin.
Naměřené hodnoty pro oba analyzované přístupy jsou shrnuty v tabulce~\ref{tab:token-usage}.

\begin{table}[H]
    \centering
    \resizebox{\textwidth}{!}{%
        \begin{tabular}{llllll}
            \toprule
            \textbf{Strategie} & \textbf{Odevzdání} & \textbf{Čas (s)} & \textbf{Problémy} & \textbf{Vstupní tokeny} & \textbf{Výstupní tokeny} \\
            \midrule
            Chain-of-thought   & 1233               & 157,09           & 3                 & 9\,584                  & 5\,550                   \\
            Chain-of-thought   & 1223               & 171,21           & 4                 & 9\,687                  & 6\,404                   \\
            Chain-of-thought   & 1224               & 107,44           & 2                 & 7\,247                  & 4\,314                   \\
            Chain-of-thought   & 1215               & 103,48           & 2                 & 9\,620                  & 3\,863                   \\
            Chain-of-thought   & 1228               & 131,69           & 4                 & 9\,126                  & 5\,614                   \\
            \midrule
            One-shot           & 1233               & 16,85            & 5                 & 2\,070                  & 590                      \\
            One-shot           & 1223               & 12,64            & 5                 & 891                     & 478                      \\
            One-shot           & 1224               & 14,02            & 5                 & 1\,218                  & 570                      \\
            One-shot           & 1215               & 13,04            & 2                 & 4\,508                  & 305                      \\
            One-shot           & 1228               & 21,88            & 6                 & 1\,527                  & 564                      \\
            \bottomrule
        \end{tabular}
    }
    \caption{Detailní výsledky experimentu s modelem Qwen3-Coder:30B na pěti studentských odevzdáních z předmětu UPR. Každé odevzdání bylo analyzováno oběma strategiemi.}
    \label{tab:cot-vs-oneshot-detail}
\end{table}

\begin{table}[H]
    \centering
    \resizebox{\textwidth}{!}{%
        \begin{tabular}{lrrrr}
            \toprule
            \textbf{Přístup} & \textbf{Vstup min--max} & \textbf{Vstup průměr} & \textbf{Výstup min--max} & \textbf{Výstup průměr} \\
            \midrule
            One-shot         & 891--4~508              & \tavg{2~043}          & 305--590                 & \tavg{501}             \\
            Chain-of-thought & 7~247--9~687            & \tavg{9~053}          & 3~863--6~404             & \tavg{5~149}           \\
            \bottomrule
        \end{tabular}
    }
    \caption{Naměřená spotřeba tokenů při analýze studentských odevzdání viz \tabref{tab:cot-vs-oneshot-detail}.}
    \label{tab:token-usage}
\end{table}

\begin{table}[H]
    \centering
    \begin{tabular}{lll}
        \toprule
        \textbf{Metrika} & \textbf{Chain-of-thought} & \textbf{One-shot} \\
        \midrule
        Průměrný čas (s)             & 134,18    & 15,69 \\
        Průměrný počet problémů      & 3,0       & 4,6 \\
        Průměrné vstupní tokeny      & 9\,053    & 2\,043 \\
        Průměrné výstupní tokeny     & 5\,149    & 501 \\
        Poměr rychlosti              & 1$\times$ & 8,5$\times$ \\
        \bottomrule
    \end{tabular}
    \caption{Průměrné hodnoty porovnání chain-of-thought a one-shot přístupu na modelu Qwen3-Coder:30B.}
    \label{tab:cot-vs-oneshot-summary}
\end{table}

Pro typickou analýzu jednoho studentského odevzdání je tedy při přístupu one-shot potřeba přibližně \tavg{2043} vstupních a \tavg{501} výstupních tokenů.
Při použití chain-of-thought přístupu, který zahrnuje více kroků analýzy, se tyto počty zvyšují na \tavg{9053} vstupních a \tavg{5149} výstupních tokenů.
Spotřeba tokenů je tedy v tomto případě přibližně čtyř až desetkrát vyšší.
Celkové náklady na jedno odevzdání jsou pak dány cenou za token konkrétního modelu.

U on-premise nasazení jsou náklady primárně jednorázové.
Zahrnují pořízení hardwaru, především grafických karet s dostatečnou kapacitou VRAM pro inferenci (viz \secref{subsec:llm-computational-cost}).
Průběžné náklady tvoří elektřina, chlazení a lidská práce spojená se správou a aktualizací modelu.
Na rozdíl od cloudového API jsou však marginální náklady na každé další zpracované odevzdání prakticky nulové, což činí on-premise řešení ekonomicky výhodnějším při velkém objemu požadavků.

\endinput